% Options for packages loaded elsewhere
\PassOptionsToPackage{unicode}{hyperref}
\PassOptionsToPackage{hyphens}{url}
%
\documentclass[
  11pt,
]{article}
\usepackage{lmodern}
\usepackage{amssymb,amsmath}
\usepackage{ifxetex,ifluatex}
\ifnum 0\ifxetex 1\fi\ifluatex 1\fi=0 % if pdftex
  \usepackage[T1]{fontenc}
  \usepackage[utf8]{inputenc}
  \usepackage{textcomp} % provide euro and other symbols
\else % if luatex or xetex
  \usepackage{unicode-math}
  \defaultfontfeatures{Scale=MatchLowercase}
  \defaultfontfeatures[\rmfamily]{Ligatures=TeX,Scale=1}
\fi
% Use upquote if available, for straight quotes in verbatim environments
\IfFileExists{upquote.sty}{\usepackage{upquote}}{}
\IfFileExists{microtype.sty}{% use microtype if available
  \usepackage[]{microtype}
  \UseMicrotypeSet[protrusion]{basicmath} % disable protrusion for tt fonts
}{}
\makeatletter
\@ifundefined{KOMAClassName}{% if non-KOMA class
  \IfFileExists{parskip.sty}{%
    \usepackage{parskip}
  }{% else
    \setlength{\parindent}{0pt}
    \setlength{\parskip}{6pt plus 2pt minus 1pt}}
}{% if KOMA class
  \KOMAoptions{parskip=half}}
\makeatother
\usepackage{xcolor}
\IfFileExists{xurl.sty}{\usepackage{xurl}}{} % add URL line breaks if available
\IfFileExists{bookmark.sty}{\usepackage{bookmark}}{\usepackage{hyperref}}
\hypersetup{
  hidelinks,
  pdfcreator={LaTeX via pandoc}}
\urlstyle{same} % disable monospaced font for URLs
\usepackage[margin=1in]{geometry}
\setlength{\emergencystretch}{3em} % prevent overfull lines
\providecommand{\tightlist}{%
  \setlength{\itemsep}{0pt}\setlength{\parskip}{0pt}}
\setcounter{secnumdepth}{-\maxdimen} % remove section numbering
\usepackage{esint}
\setmathfont{Asana Math}
\setlength{\emergencystretch}{3em}

\author{}
\date{}

\begin{document}

\hypertarget{from-weighted-coercivity-to-uniformly-local-einstein-repair}{%
\section{From weighted coercivity to uniformly local Einstein
repair}\label{from-weighted-coercivity-to-uniformly-local-einstein-repair}}

24 September 2026. This continuation removes the unnormalized global
\(L^2\) residual from the preceding
\href{20260924_local_to_global_scope.md}{diameter-free repair
criterion}, on a specified growth--coercivity regime. The additional
information is a quantitative bound on propagation through the whole
manifold, not another smallness condition on the residual's support.

\hypertarget{statement}{%
\subsection{1. Statement}\label{statement}}

For a smooth metric on a closed four-manifold use the Bianchi-gauged
operator \[
 \Phi_g(h)=\operatorname{Ric}_h+h+\delta_h^*(B_gh),\qquad
 B_gu=\delta_gu+\tfrac12d\operatorname{tr}_gu.
\] Its derivative at \(g\) is \[
 L_g=\tfrac12\Delta_{L,g}+I
     =\tfrac12\nabla^*\nabla+V_g,                            \tag{1}
\] where \(V_g\) is the symmetric curvature endomorphism. Unlike the
earlier absolute spectral gap, the condition used here is
\textbf{positive coercivity}: \[
 \int_M\langle L_gu,u\rangle\,d\mu_g
              \ge\lambda\int_M|u|^2\,d\mu_g
 \quad\text{for every smooth symmetric two-tensor }u.         \tag{2}
\] No kernel or low mode has been projected out.

\textbf{Theorem A (pointwise repair with controlled volume growth).} Fix
\(v,K_1,C_G,\lambda>0\), \(d,H\ge0\), and \(0<\alpha<1\), with \[
                              H^2<8\lambda.                 \tag{3}
\] There are \(\varepsilon,C>0\), depending only on these data, such
that the following holds. Suppose a closed connected oriented smooth
four-manifold \(M\), with \(\pi_2(M)=0\), carries a smooth metric \(g\)
satisfying \[
 \inf_x\operatorname{Vol}_gB_1(x)\ge v,\qquad
 \|\nabla\operatorname{Ric}_g\|_\infty\le K_1,                 \tag{4}
\] the positive coercivity (2), and the actual-manifold growth bound \[
 \operatorname{Vol}_g B_R(x)\le C_G(1+R)^d e^{HR}
       \quad(x\in M,\ R\ge0).                               \tag{5}
\] If \[
 \delta_-(g)=\|g^{-1}\operatorname{Ric}_g+I\|_{\infty,\mathrm{op}}
                    <\varepsilon,                          \tag{6}
\] then the same smooth manifold has a smooth metric \(h\) with \[
 \operatorname{Ric}_h=-h,\qquad
 \|h-g\|_{C^{2,\alpha}_{\mathrm{uloc}}}
                       \le C\delta_-(g)^{1-\alpha}.           \tag{7}
\] Here the uniformly local Hölder norm is the supremum of coefficient
norms in a uniformly controlled small harmonic atlas, not a sum over
that atlas. There is no diameter, total-volume, total-residual, or
support-volume bound. The theorem does not require a gauge class or a
prescribed complex structure.

The topological assumption is one supplier of local bounded geometry. If
uniform \(C^{2,\alpha}\) harmonic charts are provided directly, the same
theorem holds without \(\pi_2(M)=0\) and without using (4) to obtain
them; in that formulation the residual hypothesis is
\(\|\operatorname{Ric}+g\|_{C^{0,\alpha}_{\mathrm{uloc}}}<\varepsilon\).
Thus the main analytic conclusion is a genuine uniformly local
\(C^{0,\alpha}\to C^{2,\alpha}\) inverse and repair estimate.

\hypertarget{the-inverse-estimate-both-decay-and-growth-are-retained}{%
\subsection{2. The inverse estimate: both decay and growth are
retained}\label{the-inverse-estimate-both-decay-and-growth-are-retained}}

Choose once and for all \[
 H/2<a<\sqrt{2\lambda},\qquad c_a=\lambda-a^2/2>0.             \tag{8}
\] For a fixed centre \(x\), put \(\phi_x(y)=e^{-a d_g(x,y)}\). This
positive Lipschitz function satisfies \(|d\phi_x|\le a\phi_x\) almost
everywhere. No differentiability of the distance function at its cut
locus is assumed.

For any smooth tensor \(u\), integration by parts gives the exact
identity \[
 \int\langle L_g(\phi_xu),\phi_xu\rangle
 =\int\phi_x^2\langle L_gu,u\rangle
       +\tfrac12\int|d\phi_x|^2|u|^2.                       \tag{9}
\] For Lipschitz \(\phi_x\), use density in \(H^1\) in the
quadratic-form identity. The zeroth-order terms agree on both sides; the
factor \(1/2\) is the leading coefficient in (1).

If \(L_gu=f\), (2), (8), Cauchy--Schwarz, and (9) give \[
 c_a\|\phi_xu\|_2^2
 \le\int\phi_x^2\langle f,u\rangle
 \le\|\phi_xf\|_2\|\phi_xu\|_2,
\quad
 \|\phi_xu\|_2\le c_a^{-1}\|\phi_xf\|_2.                    \tag{10}
\] The right side can now be bounded by a \textbf{local supremum},
rather than the global \(L^2\) norm. Splitting into unit annuli and
using (5), \[
 \begin{aligned}
 \int_M e^{-2a d_g(x,y)}\,d\mu_g(y)
 &\le C_G e^H\sum_{m=0}^{\infty}
            (m+2)^d e^{-(2a-H)m}\ =:J_a<\infty.
 \end{aligned}                                             \tag{11}
\] The constant is independent of \(x\), diameter, and total volume.
Therefore \[
 \|u\|_{L^2(B_1(x))}
        \le e^a c_a^{-1}\sqrt{J_a}\,\|f\|_\infty.             \tag{12}
\]

In the regime of Theorem A, \(\delta_-<1/2\) bounds Ricci uniformly. The
local regularity theorem in the preceding continuation, using
\(\pi_2=0\) and the unit-ball lower bound, supplies uniform harmonic
radii. The Ricci derivative bound and the harmonic metric equation then
supply uniform \(C^{2,\alpha}\) metric bounds on smaller charts. In
particular the coefficients of (1), including \(V_g\), have uniform
\(C^{0,\alpha}\) bounds.

Local Moser estimates for the connection Laplacian with bounded
potential give, for a fixed \(0<\rho\le1\), \[
 |u|(x)\le C_{\rm loc}
       \bigl(\|u\|_{L^2(B_\rho(x))}+\|f\|_{C^0(B_\rho(x))}\bigr).
\] Combining this with (12), then applying local Schauder estimates,
proves \[
 \boxed{\quad
 \|u\|_{C^{2,\alpha}_{\rm uloc}}
     \le C_I\|L_gu\|_{C^{0,\alpha}_{\rm uloc}}.
 \quad}                                                     \tag{13}
\] On each closed manifold, positive coercivity and elliptic Fredholm
theory give existence and uniqueness of \(u=L_g^{-1}f\). Equation (13)
makes the inverse uniform across the family. We have not deduced (13)
from an unweighted \(L^2\) gap alone: (11) is the additional global
step.

For polynomial growth, \(H=0\), every \(\lambda>0\) permits a choice in
(8). For exponential growth, the sufficient condition (3) is the precise
comparison needed by this argument. It is not claimed optimal. In
particular failure of (3) does not imply failure of an Einstein inverse.

\hypertarget{nonlinear-repair}{%
\subsection{3. Nonlinear repair}\label{nonlinear-repair}}

In the uniformly local Hölder spaces put \[
 E=\operatorname{Ric}_g+g,\qquad
 \Phi_g(g+u)=E+L_gu+N_g(u).
\] Uniform \(C^{2,\alpha}\) bounds give, in a fixed sufficiently small
ball, \[
 \|N_g(u)-N_g(z)\|_{C^{0,\alpha}_{\rm uloc}}
 \le C_N\bigl(\|u\|_{C^{2,\alpha}_{\rm uloc}}
             +\|z\|_{C^{2,\alpha}_{\rm uloc}}\bigr)
                \|u-z\|_{C^{2,\alpha}_{\rm uloc}}.             \tag{14}
\] This is a local Hölder product estimate: neither summation over all
charts nor a volume factor appears. By (4) and tensor interpolation, \[
 \|E\|_{C^{0,\alpha}_{\rm uloc}}\le C\delta_-^{1-\alpha}.       \tag{15}
\] Equations (13)--(15) make \(u\mapsto-L_g^{-1}(E+N_g(u))\) a
contraction in the ball of radius
\(2C_I\|E\|_{C^{0,\alpha}_{\rm uloc}}\) when (6) is sufficiently small.
The resulting \(h=g+u\) is positive definite and has negative definite
Ricci tensor. Applying \(B_h\) to \(\Phi_g(h)=0\) gives, for \(b=B_gh\),
\[
              \nabla_h^*\nabla_h b-\operatorname{Ric}_h(b)=0.
\] Integration on the closed manifold gives \(b=0\), so \(h\) is
Einstein. Elliptic bootstrapping gives smoothness. This proves Theorem
A.

\hypertarget{a-nonempty-large-diameter-family-with-distributed-residual}{%
\subsection{4. A nonempty large-diameter family with distributed
residual}\label{a-nonempty-large-diameter-family-with-distributed-residual}}

Fix an oriented closed hyperbolic four-manifold \((N,h)\) with
\(\operatorname{Ric}_h=-h\) and an epimorphism
\(\varphi:\pi_1(N)\to\mathbb Z\). Such compact examples exist by
\href{https://math.umd.edu/~jjm/papers/bettinumber.pdf}{Millson's
construction}, which gives nonzero first Betti number in every dimension
at least three. Let \((\widehat N,\widehat h)\) be the cyclic infinite
cover and \((N_m,h_m)\) its quotient by \(m\mathbb Z\).

\textbf{Lemma B (the correct cover-growth bound).} There is
\(C_N<\infty\), depending only on this fixed cyclic cover, such that \[
 \operatorname{Vol}_{h_m}B_R(x)\le C_N(1+R)
                 \quad(m\ge1,\ R\ge0).                      \tag{16}
\]

\textbf{Proof.} The deck group \(\mathbb Z\) acts properly and
cocompactly on \(\widehat N\). Choose a bounded fundamental region of
finite volume. The orbit map of this action is a quasi-isometry:
equivalently, only \(O(1+R)\) translates of the region can meet a
radius-\(R\) ball, with constants depending on the fixed region and deck
generator. Counting their volumes proves the linear bound on
\(\widehat N\), uniformly in the centre. Every radius-\(R\) ball in
\(N_m\) is the image of a radius-\(R\) ball in \(\widehat N\), by
lifting shortest paths from its centre. Its volume is at most the volume
upstairs, proving (16). \(\square\)

The cover used in this argument is \textbf{not} the universal hyperbolic
cover, and (16) is not asserted for arbitrary finite covers. Since
\(\operatorname{Vol}(N_m)=m\operatorname{Vol}(N)\), (16) also shows
\(\operatorname{diam}(N_m)\to\infty\).

Choose \(\eta>0\) small and consider any smooth metrics \(g_m\) with \[
                  \|g_m-h_m\|_{C^3_{\rm uloc}(h_m)}\le\eta.   \tag{17}
\] The constants below are uniform in \(m\). The metrics are uniformly
bilipschitz to \(h_m\), so (16) remains a polynomial bound with adjusted
constant, and unit balls stay uniformly noncollapsed. They have
\(\pi_2=0\), bounded Ricci derivative, and sectional curvature uniformly
close to \(-1/3\).

For the exact \(h_m\), Koiso's estimate gives \[
                \int\langle L_{h_m}u,u\rangle
                              \ge\tfrac13\|u\|_2^2.          \tag{18}
\] Its almost-Einstein version gives, for \(\eta\) sufficiently small,
positive coercivity with \(\lambda=1/4\) for every \(g_m\) in (17). Thus
\(H=0,d=1,\lambda=1/4,a=1/2\) work simultaneously, with \[
 c_a=\tfrac18,\qquad
 J_a\le C_N\sum_{j=0}^{\infty}(j+2)e^{-j}
        =C_N\frac{2-e^{-1}}{(1-e^{-1})^2}.                    \tag{19}
\] Theorem A applies to every metric in (17) whose pointwise Ricci
residual is below its fixed threshold. The required threshold is
independent of the number of cyclic cells, and the metrics need not be
deck-invariant.

This family contains residuals which are small everywhere but not small
in global \(L^2\). For an explicit example, choose a ball in one fixed
fundamental cell and a smooth bump \(\beta\) supported there and equal
to one on a smaller ball. On its \(m\) disjoint translates choose
arbitrary coefficients \(c_j\in[1,2]\), and put \[
 f_m=t\sum_{j=1}^m c_j\beta_j,\qquad g_m=e^{2f_m}h_m.          \tag{20}
\] For small fixed \(t>0\), (17), all uniform constants and
\(\delta_-(g_m)\le C t\) hold. On each smaller ball the conformal factor
is constant, so \(\operatorname{Ric}_{g_m}+g_m =(e^{2tc_j}-1)h_m\).
Consequently \[
 \|\operatorname{Ric}_{g_m}+g_m\|_{L^2(g_m)}
                           \ge c t\sqrt m.                  \tag{21}
\] For sufficiently small fixed \(t\), Theorem A repairs every \(g_m\),
although (21) diverges. This is a real gain over the earlier mixed-norm
criterion, not merely a change from one notation for the residual to
another. The underlying covers already admit hyperbolic metrics; this
application establishes a \textbf{uniform repair neighbourhood for
distributed perturbations}, not previously unknown Einstein topologies.

\hypertarget{scope-relative-to-classical-analysis-and-the-yau26-program}{%
\subsection{5. Scope relative to classical analysis and the Yau26
program}\label{scope-relative-to-classical-analysis-and-the-yau26-program}}

The weighted quadratic-form identity is classical; see
\href{https://arxiv.org/html/2502.16692v2}{Jäckel, Remark 2.6}. His
Lemma 2.4 supplies the almost-Einstein Koiso estimate used in (18),
after scaling from \(\operatorname{Ric}=-(n-1)g\) to
\(\operatorname{Ric}=-g\). The local Moser estimate is his Lemma 2.2.
The new composition here uses the \textbf{actual quotient's slow volume
growth} to convert weighted residual norms into uniformly local norms,
and then uses the preceding topological regularity result as a source of
controlled charts.

Jäckel's Theorem 1.1 permits collapse, under its negative-sectional and
hyperbolic-thin-part hypotheses, but retains a small global \(L^2\)
residual. The present argument requires local noncollapse and (3),(5),
but permits errors spread throughout arbitrarily large volume.
\href{https://numdam.org/item/10.5802/tsg.367.pdf}{Premoselli's repair
analysis} uses a different weighted inverse tailored to its gluing
regions and source support. None of these results is an unrestricted
replacement for the others. No global-priority claim is made for the
weighted principle.

There are two important mathematical exclusions from the present
theorem. First, the earlier assumption \(\|L_gu\|_2\ge\lambda\|u\|_2\)
is not interchangeable with (2); negative eigenvalues invalidate the
weighted coercivity step. Second, products of hyperbolic curves have
genuine Teichmüller kernels, so they do not satisfy (2). An obstruction
equation cannot be bypassed by deleting those tensors. The growth
comparison also concerns all balls on the actual manifold; a global
entropy number without uniform finite-radius bounds does not by itself
supply (11).

The scientific increment is therefore specific: the prior no-diameter
criterion required \(\delta^{1-\alpha}+\|E\|_2\) to be small; here a
verified growth--coercivity condition makes \textbf{pointwise small
Ricci error} sufficient in a family of unbounded diameter and volume. It
remains a restricted sufficient answer to the original
same-smooth-manifold existence question, not a dimension-only Yau26
threshold.

\hypertarget{reproduction-scope}{%
\subsection{6. Reproduction scope}\label{reproduction-scope}}

\texttt{check\_uniformly\_local\_repair.py} verifies the factor in (9),
the growth threshold and constants in (19), and the four-dimensional
scaling in (21). These are finite algebraic checks. They do not certify
the elliptic, covering-space or nonlinear arguments above.

\end{document}
